%----------------------------------------------------------------------------------------------------
% START OF CHAPTER
%----------------------------------------------------------------------------------------------------
\chapter{May 18}
\paragraph{Topics}
\textit{Pronouns recap, Conversation}

\section{Grammar} \label{0518-Grammar}
\subsection{Pronouns recap}

\begin{itemize}
  \item \textbf{Subject pronouns}. Who is doing the action of the sentence: I, you, we, you (pl.), them $\rightarrow$ \iti{io, tu, noi, voi, loro}
  \item \textbf{Direct object pronouns}: What/whom is being acted on: me, you, him/it, her/it, us, you (pl.), them) $\rightarrow$ \iti{mi, ti, lo, la, ci, vi, li, le}
  \item \textbf{Indirect object pronouns}: Who is the recipient of that action: to/for me, you, him, her, us, you (pl.), them $\rightarrow$ \iti{mi, ti, gli, le, ci, vi, gli}
  \item \textbf{Disjunctive pronouns}: Extra info with prepositions like `for', `between', etc: <preposition> + me, you, him, her, us, you (pl.), them $\rightarrow$ \iti{me, te, lui, lei, noi, voi, loro}
\end{itemize}

\subsubsection{Example}
Consider: `Luigi gave the book to Luca yesterday, not to Isabella.'
With pronouns this becomes: `He gave it to him yesterday, not to her.'


\begin{itemize}[nosep]
  \item{Luigi $\rightarrow$ subject pronoun `He' $\rightarrow$ \iti{lui}}
  \item{gave $\rightarrow$ verb \iti{ha dato}}
  \item{the book $\rightarrow$ direct object pronoun `it' $\rightarrow$ \iti{lo}}
  \item{to Luca $\rightarrow$ indirect object pronoun `to him' $\rightarrow$ \iti{gli}}
  \item{yesterday $\rightarrow$ \iti{ieri}}
  \item{not $\rightarrow$ \iti{non}}
  \item{to Isabella $\rightarrow$ preposition `to' and disjunctive pronoun `her' $\rightarrow$ \iti{a lei}}
\end{itemize}

\paragraph{Step by step walkthrough}
\begin{enumerate}[nosep]
  \item{Focus on the main clause, `He gave it to him'.}
  \item{Rule: indirect pronoun comes before direct, so we translate right-to-left: [to him] [it] [he gave].}
    \footnote{Exception to this rule is when \iti{loro} when used as indirect pronoun - it comes after the verb.}
  \item{This gets us: \iti{Gli} + \iti{lo} + \iti{ha dato}.}
  \item{Rule: indirect + direct combos may contract or change form - see \ref{tab:tIndirectDirectCombinations}.}
  \item{New form: \iti{Glielo} + \iti{ha dato}.}
  \item{Add in `yesterday' and the preposition `not' and the disjunctive `to her'.}
  \item{\iti{Glielo ha dato} + \iti{ieri} + \iti{non} + \iti{a lei}.}
  \item{Result: \iti{Glielo ha dato ieri non a lei.}}
\end{enumerate}


When indirect and direct pronouns are used together, their forms change to flow more naturally.
For example, \textit{mi li} becomes \iti{me li}. See \ref{tab:tIndirectDirectCombinations}.

\begin{table}[ht]
\centering
    \begin{tabular}{ c c | l | l r }
      \textbf{indirect} & \textbf{+ direct} & \textbf{$\rightarrow$ Result} & \textit{Example} &  \\
         \hline
         \iti{mi}       & \iti{lo/la/li/le} & \iti{me lo/la/li/le}        & \iti{Me lo dice}  & He says it to me. \\
         \iti{ti}       & \iti{lo/la/li/le} & \iti{te lo/la/li/le}        & \iti{Te lo porto} & I bring it to you. \\
         \\
         \iti{gli}      & \iti{lo/la/li/le} & \iti{glielo/la/li/le}       & \iti{Glieli lo già dati} & I already gave them to him. \\
         \iti{le}       & \iti{lo/la/li/le} & \iti{glielo/la/li/le}       & \iti{Glielo scrivo} & I write it to her. \\
         \\
         \iti{ci}       & \iti{lo/la/li/le} & \iti{ce lo/la/li/le}        & \iti{Ce li danno}  & They give them to us. \\
         \iti{vi}       & \iti{lo/la/li/le} & \iti{ve lo/la/li/le}        & \iti{Ve li mostro} & I show them to you (pl.) \\
         \\
         \iti{loro}     & \iti{lo/la/li/le} & \iti{lo/la/li/le} + verb + \iti{loro}  & \iti{Lo do loro} & I give it to them.
    \end{tabular}
    \caption{Indirect and Direct Object pronoun combinations}
    \label{tab:tIndirectDirectCombinations}
\end{table}

\subsection{Pronouns after the verb}
\begin{mdframed}
  Normally pronouns go before the verb: \iti{Me \sout{lo} dice}, for example. However, with modal verbs (can, must, want),
  imperatives, and some expressions (`without', `instead of', `after', + verb) they go after.
\end{mdframed}

\paragraph{Examples}
\begin{itemize}
  \item{I can buy it} \\
    \iti{Posso comprar\sout{lo}.}

  \item{Eat it! (gelato)} \\
  \iti{Mangialo!}

  \item{Instead of reading it, buy it.} \\
  \iti{Invece di leggerlo, comprarlo.}

  \item{I can write it, without studying it.} \\
  \iti{Posso scriverlo, senza studiarlo.}

  \item{After driving it, I don't want to buy it.} \\
  \iti{Dopo guidarla, non voglio comprarla.}

  \item{I can't buy it without trying it.} \\
  \iti{Non posso comprarlo senza proverlo.}

  \item{I don't want to buy you gelato now.} \\
  \iti{Non voglio comprarti gelato adesso.}

  \item{I don't want to buy her gelato now.} \\
  \iti{Non voglio comprarle gelato adesso.}

\end{itemize}


\section{Conversation} \label{0518-Conversation}
\begin{itemize}
  \item [P] \iti{Sono stanca perché sto lavorando troppo}
  \item [A] \iti{Mi dispiace. Tu dovresti parlare con il tuo capo.}
  \item [F] \iti{Buona idea. Non pensare a questo. Comprati un gelato.}
  \item [S] \iti{Non ti preoccupare. Andiamo al ristorante questa sera.}
  \item [P] \iti{Grazie mille, ma, io devo andare al bar con una mia amica.}
  \item [A] \iti{Facciamo domani sera.}
  \item [F] \iti{Io sono disponibile domani sera verso le sette di sera.}
  \item [P] \iti{Perfetto. Ci vediamo domani sera.}
\end{itemize}


%----------------------------------------------------------------------------------------------------
% END OF CHAPTER
%----------------------------------------------------------------------------------------------------
